```latex
\documentclass[a4paper,10pt]{article}

\usepackage[utf8]{inputenc}
\usepackage[T1]{fontenc}
\usepackage{lmodern}
\usepackage[french]{babel}
\usepackage{geometry}
\usepackage{xcolor}
\usepackage{hyperref}
\usepackage{titlesec}
\usepackage{fontawesome5}
\usepackage{enumitem}

% --- Configuration de la page ---
\geometry{
    left=1.5cm,
    right=1.5cm,
    top=1.5cm,
    bottom=1.5cm
}
\pagestyle{empty}

% --- Couleurs et polices ---
\definecolor{primary}{HTML}{005A9C} % Un bleu professionnel
\definecolor{secondary}{HTML}{4F4F4F} % Gris foncé pour le texte
\definecolor{lightgray}{HTML}{E0E0E0} % Gris clair pour les séparateurs

\usepackage{helvet}
\renewcommand{\familydefault}{\sfdefault}
\color{secondary}

% --- Liens cliquables ---
\hypersetup{
    colorlinks=true,
    urlcolor=primary,
    linkcolor=primary,
}

% --- Personnalisation des sections ---
\titleformat{\section}
  {\Large\bfseries\color{primary}}
  {}
  {0em}
  {}[\titlerule\vspace{-0.5ex}]
\titlespacing*{\section}{0pt}{1.5em}{1em}

% --- Informations personnelles ---
\newcommand{\candidatename}{Yassine Dridi}
\newcommand{\candidateemail}{yassine.dridi@dauphine.tn}
\newcommand{\candidatelocation}{Tunis, Tunisie}
\newcommand{\candidatelinkedin}{linkedin.com/in/yassinedridi} % Placeholder

% --- Document ---
\begin{document}

% ====================
% EN-TÊTE
% ====================
\begin{center}
    {\Huge\bfseries\scshape\candidatename}
    \vspace{0.5em}

    \begin{tabular}{c c c}
        \faEnvelope & \faMapMarkerAlt & \faLinkedin \\
        \href{mailto:\candidateemail}{\candidateemail} & \candidatelocation & \href{https://\candidatelinkedin}{\candidatelinkedin}
    \end{tabular}
\end{center}

% ====================
% PROFIL
% ====================
\section*{Profil}
Data Scientist avec une solide expérience en Machine Learning, analyse de données et modélisation statistique. Passionné par la transformation de données brutes en informations stratégiques, je cherche à appliquer mes compétences en Python, SQL et visualisation de données pour résoudre des problématiques complexes au sein de Modulaire Group en tant que Data Scientist Intern.

% ====================
% COMPÉTENCES
% ====================
\section*{Compétences Techniques}
\begin{itemize}[leftmargin=*, label={}, itemsep=0.2em]
    \item \textbf{Programmation \& Analyse de Données :} \textbf{Python} (Pandas, NumPy, Scikit-learn), SQL
    \item \textbf{Machine Learning :} Modèles de régression et classification, Clustering, Feature Engineering, Validation de modèles
    \item \textbf{Modélisation Statistique :} Tests d'hypothèses, Analyse de séries temporelles, Probabilités
    \item \textbf{Visualisation de Données :} Matplotlib, Seaborn, Plotly (Connaissances de Tableau)
    \item \textbf{Bases de Données :} PostgreSQL, MySQL
\end{itemize}

% ====================
% EXPÉRIENCE PROFESSIONNELLE
% ====================
\section*{Expérience Professionnelle}

\begin{tabular*}{\textwidth}{@{}l@{\extracolsep{\fill}}r@{}}
\textbf{Data Scientist} & \textit{Janv. 2021 -- Présent} \\
\textit{Tech Solutions Inc.} & \textit{Tunis, Tunisie} \\
\end{tabular*}
\begin{itemize}[leftmargin=*, topsep=0.3em, itemsep=0.2em]
    \item Développement et déploiement de modèles de \textbf{Machine Learning} pour prédire le comportement client, résultant en une amélioration de 15\% de la rétention.
    \item \textbf{Analyse de données} exploratoire sur de grands jeux de données (>10M d'entrées) avec \textbf{Python} et \textbf{SQL} pour identifier des tendances et des opportunités de croissance.
    \item Création de dashboards de \textbf{visualisation de données} pour communiquer les résultats des analyses aux équipes marketing et produit.
    \item Optimisation des requêtes \textbf{SQL} pour le traitement et la manipulation des données, réduisant les temps d'exécution de 30\%.
\end{itemize}

\vspace{1em}

\begin{tabular*}{\textwidth}{@{}l@{\extracolsep{\fill}}r@{}}
\textbf{Data Analyst} & \textit{Juin 2019 -- Déc. 2020} \\
\textit{Data Insights Corp.} & \textit{Tunis, Tunisie} \\
\end{tabular*}
\begin{itemize}[leftmargin=*, topsep=0.3em, itemsep=0.2em]
    \item Extraction et nettoyage de données provenant de sources multiples via des scripts Python et des requêtes SQL.
    \item Production de rapports hebdomadaires et mensuels sur les indicateurs de performance clés (KPIs).
    \item Collaboration avec les équipes métiers pour définir les besoins en analyse et traduire les questions business en requêtes de données.
\end{itemize}

% ====================
% PROJETS ACADÉMIQUES ET PERSONNELS
% ====================
\section*{Projets Pertinents}

\begin{tabular*}{\textwidth}{@{}l@{\extracolsep{\fill}}r@{}}
\textbf{Modèle de Prédiction de Churn Client} & \textit{Projet Universitaire} \\
\end{tabular*}
\begin{itemize}[leftmargin=*, topsep=0.3em, itemsep=0.2em]
    \item Conception d'un pipeline complet de \textbf{Machine Learning} pour un dataset de télécommunications, de la préparation des données à l'évaluation du modèle (Logistic Regression, Random Forest).
    \item Utilisation de techniques de \textbf{Data Visualization} pour présenter les facteurs d'influence principaux du churn.
    \item \textbf{Technologies :} Python, Scikit-learn, Pandas, Matplotlib, Seaborn.
\end{itemize}

% ====================
% FORMATION
% ====================
\section*{Formation}
\begin{tabular*}{\textwidth}{@{}l@{\extracolsep{\fill}}r@{}}
\textbf{Master en Mathématiques et Informatique Appliquées (MIA)} & \textit{2017 -- 2019} \\
\multicolumn{2}{@{}l}{\textit{Université Paris-Dauphine | Tunis}} \\
\end{tabular*}
\textit{Spécialisation en Data Science et Analyse Statistique.}

% ====================
% COMPÉTENCES TRANSVERSES
% ====================
\section*{Compétences Transverses}
\begin{itemize}[leftmargin=*, label={}, itemsep=0.2em, topsep=0.3em]
    \item \textbf{Communication :} Présentation claire des résultats techniques à des audiences non-techniques.
    \item \textbf{Résolution de problèmes :} Approche analytique et structurée pour décomposer des problèmes complexes.
    \item \textbf{Travail d'équipe :} Expérience de collaboration en mode agile avec des équipes pluridisciplinaires.
\end{itemize}

\end{document}
```